\chapter{Scales, Strings, and Twenty-Six Dimensions}

\lectureoverview{What sets the spacing of a string spectrum? The answer is
the tension, a quantity with dimensions of mass squared. For fundamental
strings that observation leads toward Planck units, the practical
impossibility of direct excitation, and the onset of strong gravity. After a
deliberate break, we return to the bosonic open-string ground state.
Regulating its zero-point energy gives the historical
\emph{easy} route to twenty-four transverse directions and hence twenty-six
spacetime dimensions.}

\section{The Scale We Had Left Out}

So far we have spoken about vibrating and rotating strings without saying
what one unit of excitation means. A hadron can be pictured as a string, and
its oscillators can be excited, but how large is the jump from one state to
the next? What determines it?

There is only one dimensional parameter in the elementary string model. To
find it, begin with an ordinary stretched spring. If its endpoints are
separated by an amount \(\ell\), Hooke's law gives
\begin{equation}
  E_{\mathrm{spring}}=\frac{k\ell^2}{2}.
  \label{eq:l05-hooke}
\end{equation}
The factor of one-half was supplied by the class after the first version was
written. Nothing essential will depend on such numerical factors in this
opening estimate.

\begin{figure}[H]
  \centering
  \includegraphics[width=0.72\textwidth]{lecture_05_figure_01.png}
  \caption{The stretched segment and its Hooke-law potential energy.}
  \lectureframe{00:02:28}
  \label{fig:l05-hooke-board}
\end{figure}

The energy in \eqref{eq:l05-hooke} is nonrelativistic. In the
infinite-momentum description used in the preceding lectures, internal
nonrelativistic energy is proportional not to \(m\), but to \(m^2\). We
therefore write schematically
\begin{equation}
  m^2\sim k\ell^2,
  \qquad
  m\sim \sqrt{k}\,\ell.
  \label{eq:l05-mass-length}
\end{equation}
In the string's rest frame its mass is its energy, so the second relation has
the form
\begin{equation}
  E_{\mathrm{rest}}\sim T_s\ell,
  \qquad
  T_s\equiv\sqrt{k}.
  \label{eq:l05-tension}
\end{equation}
The coefficient \(T_s\) is the string tension: the energy stored per unit
length.

\begin{figure}[H]
  \centering
  \includegraphics[width=0.78\textwidth]{lecture_05_figure_02.png}
  \caption{The board passes from \(m^2\sim k\ell^2\) to a mass proportional
  to length.}
  \lectureframe{00:05:42}
  \label{fig:l05-mass-length-board}
\end{figure}

This notation is adapted to the lecture's spring analogy. In standard
relativistic string notation one usually writes
\begin{equation}
  T_s=\frac{1}{2\pi\alpha'},
  \label{eq:l05-alpha-prime}
\end{equation}
so increments of string level produce increments of \(m^2\) of order
\(1/\alpha'\), equivalently of order the tension.
\editorialnote{The lecture's effective \(k\) is chosen so that
\(\sqrt{k}\) plays the role of tension. The convention-independent statement
is that \(T_s\) has dimensions of mass squared and that the Regge slope
satisfies \(\alpha'\propto T_s^{-1}\).}

\subsection{Dimensions and oscillator gaps}

Set \(c=\hbar=1\). Energy then has dimensions of inverse length:
\begin{equation}
  [E]=L^{-1}.
\end{equation}
Since tension is energy per unit length,
\begin{equation}
  [T_s]=L^{-2}=E^2.
  \label{eq:l05-tension-dimension}
\end{equation}
That is why a string spectrum is naturally spaced in mass squared.

\begin{figure}[H]
  \centering
  \includegraphics[width=0.68\textwidth]{lecture_05_figure_03.png}
  \caption{Natural units on the board: \(c=\hbar=1\), energy is inverse
  length, and the oscillator frequency is controlled by the spring scale.}
  \lectureframe{00:06:05}
  \label{fig:l05-natural-units-board}
\end{figure}

The ordinary oscillator makes the same point from another direction. For a
spring with mass \(\mu\),
\begin{equation}
  \omega=\sqrt{\frac{k}{\mu}},
  \qquad
  \Delta E=\hbar\omega.
\end{equation}
A stiffer spring has a higher frequency and a larger quantum gap. A string
contains infinitely many such normal modes, but their common scale is fixed
by \(T_s\). Once the tension is known, the spacing of the string spectrum is
known.

Tension also has the dimensions of force, because force times distance is
energy. Imagine anchoring one endpoint and suspending a weight from the
other. An ideal relativistic string with \(E=T_s\ell\) can support a force
set by \(T_s\), independent of its length. The hadronic tension inferred from
observed excitation energies corresponds, in this deliberately fanciful
comparison, to the weight of a truck. No truck has ever been hung from a
meson; the measured spectrum gives the tension, and the tension gives the
mechanical comparison.

Fundamental strings would be far stiffer. If they underlie electrons,
photons, and gravitons, their first oscillator excitations must lie far above
observed particle energies. In the same fanciful gravitational experiment
their tension could support something of galactic mass. The point of the
comparison is not engineering. It is that an enormous tension means both a
large excitation energy and a very small ground-state fluctuation.

\section{Units Made from Universal Constants}

What scale is enormous enough? The natural benchmark is the Planck scale.
Rather than quote it, let us construct it.

Meters, kilograms, and seconds are excellent laboratory units, but their
origins are human-sized. A meter is convenient for cloth or rope; a foot is
named after a body part; a second is a manageable fraction of a day. One
could set the proton radius equal to one and simplify some nuclear formulas,
but the proton is only one particular composite particle. Such a choice
would not simplify all of physics.

Three dimensional constants have a more universal reach:
\begin{itemize}
  \item \(c\) sets the limiting speed for every material system;
  \item \(\hbar\) appears in the uncertainty relations of every quantum
  system;
  \item \(G\) couples gravity to every form of mass-energy.
\end{itemize}
Choosing units in which \(c=\hbar=G=1\) is equivalent to constructing units
of length, time, and mass from these constants.

Their ordinary dimensions are
\begin{equation}
  \begin{aligned}
    [c]&=LT^{-1},
    & [\hbar]&=ML^2T^{-1},\\
    [G]&=L^3M^{-1}T^{-2}.
  \end{aligned}
  \label{eq:l05-dimensions-constants}
\end{equation}
The middle relation follows directly from
\(\Delta x\,\Delta p\gtrsim\hbar\). The last follows, for example, from
\begin{equation}
  a=\frac{GM}{r^2}.
\end{equation}

\begin{figure}[H]
  \centering
  \begin{minipage}[t]{0.48\textwidth}
    \centering
    \includegraphics[width=\textwidth]{lecture_05_figure_04.png}
    \par\smallskip\footnotesize (a) Units of \(c\) and \(\hbar\).
  \end{minipage}\hfill
  \begin{minipage}[t]{0.48\textwidth}
    \centering
    \includegraphics[width=\textwidth]{lecture_05_figure_05.png}
    \par\smallskip\footnotesize (b) Units of \(G\) from Newtonian
    acceleration.
  \end{minipage}
  \caption{The dimensional-analysis ingredients assembled in class.}
  \lectureframe{00:25:43--00:26:40}
  \label{fig:l05-dimensional-analysis-board}
\end{figure}

\subsection{The Planck area first}

Seek powers \(p,q,r\) for which
\begin{equation}
  G^p\hbar^q c^r
\end{equation}
has dimensions \(L^2\). Choosing area first avoids an early square root and
also anticipates the special role that area will later play in gravity. From
\eqref{eq:l05-dimensions-constants}, cancellation of mass and time gives
\begin{equation}
  -p+q=0,
  \qquad
  -2p-q-r=0.
\end{equation}
The remaining length power must satisfy
\begin{equation}
  3p+2q+r=2.
\end{equation}
Thus \(q=p\), \(r=-3p\), and \(p=1\). The Planck area is
\begin{equation}
  \ell_P^2=\frac{G\hbar}{c^3},
  \label{eq:l05-planck-area}
\end{equation}
and consequently
\begin{equation}
  \ell_P=\sqrt{\frac{G\hbar}{c^3}},
  \qquad
  t_P=\frac{\ell_P}{c},
  \qquad
  m_P=\sqrt{\frac{\hbar c}{G}}.
  \label{eq:l05-planck-units}
\end{equation}

\begin{figure}[H]
  \centering
  \includegraphics[width=0.78\textwidth]{lecture_05_figure_06.png}
  \caption{The completed dimensional argument for the Planck area and
  length.}
  \lectureframe{00:31:34}
  \label{fig:l05-planck-area-board}
\end{figure}

At order-of-magnitude precision,
\begin{equation}
  \begin{aligned}
    \ell_P&\sim10^{-35}\,\mathrm{m},
    &t_P&\sim10^{-43}\,\mathrm{s},\\
    m_P&\sim10^{-8}\,\mathrm{kg},
    &m_Pc^2&\sim10^{19}\,\mathrm{GeV}.
  \end{aligned}
  \label{eq:l05-planck-numbers}
\end{equation}

\begin{figure}[H]
  \centering
  \includegraphics[width=0.72\textwidth]{lecture_05_figure_07.png}
  \caption{The three Planck scales recorded to order of magnitude.}
  \lectureframe{00:35:12}
  \label{fig:l05-planck-values-board}
\end{figure}

The length and time are unimaginably small, while the mass is surprisingly
macroscopic by particle standards, roughly the mass of a fine dust grain.
Its rest energy is of order a few gigajoules, comparable to the chemical
energy in a tank of gasoline. The observable universe, by contrast, spans
roughly \(10^{60}\) Planck lengths and has lived for roughly
\(10^{60}\) Planck times. These comparisons are meant to give the units
physical weight, not to improve their numerical precision.

If the characteristic string length \(\ell_s=\sqrt{\alpha'}\) is near
\(\ell_P\), then the string's quantum size is near \(\ell_P\), its
oscillation period is near \(t_P\), and its first massive excitation is near
\(m_P\). The lecture allows factors of ten, one hundred, or more; the
identification is an expectation, not a dimensional identity.

\begin{classroomqa}
\audiencequestion{Is the Planck-scale mass the increase above the particle's
ground state? \lecturetimestamp{00:38:37}}
\lecturerresponse{The evenly spaced quantity is mass squared. If the ground
state is light, the first string excitation has \(m^2\) larger by an amount
of order \(m_P^2\), so its mass and required center-of-mass energy are of
order \(m_P\).}
\end{classroomqa}

\section{How Close Should the Two Scales Be?}

\begin{classroomqa}
\audiencequestion{Dimensional analysis gives \(\ell_P\), but why should that
number be the length of a string? \lecturetimestamp{00:40:21}}
\lecturerresponse{It is a physical guess supported by gravity and by the
absence of observed string structure, not a consequence of dimensional
analysis alone. The string length may be larger than the Planck length,
which means its energy scale is lower. The lecture argues only that it should
not be an ordinary particle-physics length and that pushing it far below the
Planck length runs into gravitational collapse.}
\end{classroomqa}

The relation between the two scales is in fact coupling dependent. In a
complete compactification, \(\ell_P\), \(\ell_s\), the string coupling, and
the compactification volume enter together. There is no universal equation
\(\ell_s=\ell_P\). What survives from the lecture is the scale argument:
string effects must enter very high, and sub-Planckian localization cannot
be discussed while neglecting gravity.
\editorialnote{Modern string constructions need not identify the string
length exactly with the four-dimensional Planck length. Their relation
depends on the string coupling and compactification data.}

\begin{classroomqa}
\audiencequestion{If the Planck length is treated as a smallest useful
length, is the Planck mass also the smallest possible
mass? \lecturetimestamp{00:43:27}}
\lecturerresponse{No. Ordinary particles are far lighter. The Planck mass is
the scale at which a compact object can have a gravitational radius
comparable to its quantum wavelength. In the lecture's heuristic language,
it is the scale of the lightest semiclassical black hole, with a lifetime of
order a Planck time; the last statement should not be mistaken for a precise
description of a fully quantum object.}
\end{classroomqa}

\subsection{What running couplings can tell us}

There is indirect evidence that conventional quantum field theory remains
valid to scales far beyond direct accelerator reach. The three
standard-model gauge couplings are measured at accessible energies, and
their scale dependence can be extrapolated. In the lecture's schematic
plot, the \(U(1)\), \(SU(2)\), and \(SU(3)\) couplings approach one another at
a high unification scale, perhaps a few orders of magnitude below the Planck
energy.

\begin{figure}[H]
  \centering
  \includegraphics[width=0.70\textwidth]{lecture_05_figure_08.png}
  \caption{The schematic running-coupling plot. The three gauge curves
  approach one another before the rising gravitational strength becomes
  order one.}
  \lectureframe{00:47:18}
  \label{fig:l05-running-couplings-board}
\end{figure}

This is an extrapolation, not a direct experiment at the crossing point.
Additional particles or thresholds between laboratory and unification
energies can change the slopes. The cautious conclusion is only that known
field theory shows no sign of string-sized substructure at ordinary
distances, and that a plausible string scale lies somewhere in the enormous
interval between the inferred unification and Planck scales.

\begin{classroomqa}
\audiencequestion{Which three couplings are being
extrapolated? \lecturetimestamp{00:48:28}}
\lecturerresponse{They are the gauge couplings associated with
\(U(1)\), \(SU(2)\), and \(SU(3)\): the hypercharge, weak, and strong
interactions. Electromagnetism emerges from electroweak mixing rather than
being an independent fourth gauge factor in this accounting.}
\end{classroomqa}

\begin{classroomqa}
\audiencequestion{Does the gravitational constant \(G\) itself increase as
we probe higher energy? \lecturetimestamp{00:49:30}}
\lecturerresponse{\(G\) is dimensional, so it is not by itself the correct
quantity to compare with a dimensionless gauge coupling. At a characteristic
energy \(M\), the dimensionless gravitational strength behaves
schematically as \(GM^2/(\hbar c)\). That combination grows with energy and
becomes order one at the Planck scale.}
\end{classroomqa}

The force laws make the comparison visible:
\begin{equation}
  F_{\mathrm{em}}\sim\frac{e^2}{r^2},
  \qquad
  F_{\mathrm{grav}}\sim\frac{GM^2}{r^2}.
  \label{eq:l05-force-laws}
\end{equation}
Thus \(e^2\) is compared not with \(G\), but with \(GM^2\) in natural units.

\begin{figure}[H]
  \centering
  \includegraphics[width=0.74\textwidth]{lecture_05_figure_09.png}
  \caption{Electromagnetic and gravitational force laws used to identify the
  dimensionless gravitational strength.}
  \lectureframe{00:50:50}
  \label{fig:l05-force-laws-board}
\end{figure}

\section{Why Direct String Spectroscopy Is Remote}

How would we establish experimentally that an electron is a string? We
would do what hadron spectroscopy did: strike it hard enough to excite its
internal modes and look for a tower whose masses follow the string pattern.
If the first massive level lies near \(10^{19}\,\mathrm{GeV}\), that is far
beyond any direct machine.

The total energy is not the only obstacle. It must be concentrated within a
tiny impact parameter. That brings in black holes.

\begin{classroomqa}
\audiencequestion{At these energies would an accelerator begin producing
small black holes? \lecturetimestamp{00:50:57}}
\lecturerresponse{That is the heuristic expectation once center-of-mass
energy is concentrated inside its gravitational radius. A stellar black
hole is made by gravitational potential energy squeezing a massive star. A
microscopic one would instead be made by kinetic energy focused into a tiny
collision region. In four-dimensional gravity the required scale is
Planckian.}
\end{classroomqa}

For perspective, a linear accelerator's energy gain is roughly proportional
to its length. The lecture takes about two miles for roughly
\(100\,\mathrm{GeV}\) as a simple SLAC-scale yardstick. Reaching
\(10^{19}\,\mathrm{GeV}\) at the same gradient would multiply the length by
roughly \(10^{17}\), producing a machine of galactic scale.

Even such a machine would need luminosity. Beams are not aimed with
Planck-length precision one pair at a time; enormous numbers of particles
cross so that a few collide at the desired impact parameter. Multiplying a
Planck energy by an enormous collision rate produces an equally absurd
power requirement. The estimates are intentionally rough, but their
conclusion is stable: direct string spectroscopy at the Planck scale is not
an experiment one can plausibly design.

Indirect evidence would have to be more structural: a theory that yields the
observed low-energy particles and interactions with enough specificity to be
tested, or effects of compactification that survive far below the string
scale. The lecture does not claim that such a decisive construction is
already in hand.

\subsection{Extra dimensions and the localization limit}

\begin{classroomqa}
\audiencequestion{Should the extra dimensions also have sizes near the
Planck length? \lecturetimestamp{00:56:23}}
\lecturerresponse{They are expected to be extremely small in this elementary
picture, though not fixed exactly by the argument. If a compact dimension
were large enough to explore at lower energy, new Kaluza--Klein modes and
higher-dimensional behavior would already have altered the field theory we
observe. Its successful extrapolation therefore places an upper bound on
such dimensions.}
\end{classroomqa}

\begin{classroomqa}
\audiencequestion{Why could a string or compact dimension not be many orders
of magnitude smaller than \(\ell_P\)? \lecturetimestamp{00:57:48}}
\lecturerresponse{Resolving a shorter distance requires a larger momentum.
Beyond the Planck scale the energy needed for localization has a
gravitational radius comparable to, and then larger than, the region being
probed. The experiment makes a black hole before it resolves a smaller
ordinary string.}
\end{classroomqa}

The uncertainty principle gives the estimate. Localizing to \(\Delta x\)
requires

\begin{classroomqa}
\audiencequestion{Once Planck length and Planck time are chosen, how is the
Planck mass physically determined? \lecturetimestamp{00:59:04}}
\lecturerresponse{Use either \(\hbar\) or \(G\) to convert the length scale
into an energy. The uncertainty principle says that localizing two particles
within \(\ell_P\) requires momentum of order \(\hbar/\ell_P\), hence energy
of order \(m_Pc^2\). Equivalently, it is the mass for which quantum and
gravitational lengths become comparable.}
\end{classroomqa}

\begin{equation}
  \Delta p\gtrsim\frac{\hbar}{\Delta x},
  \qquad
  E\sim c\,\Delta p.
  \label{eq:l05-localization}
\end{equation}
At \(\Delta x=\ell_P\), this energy is
\begin{equation}
  E\sim\frac{\hbar c}{\ell_P}=m_Pc^2.
\end{equation}
At the same time the Schwarzschild radius
\begin{equation}
  r_s\sim\frac{GE}{c^4}
\end{equation}
is of order \(\ell_P\), up to numerical constants. Quantum localization and
gravitational collapse meet at the same scale.

\begin{figure}[H]
  \centering
  \includegraphics[width=0.74\textwidth]{lecture_05_figure_10.png}
  \caption{The Planck length and mass formulas revisited during the
  localization argument.}
  \lectureframe{01:00:25}
  \label{fig:l05-localization-board}
\end{figure}

This completes the scale-setting argument.

\section{The Ground State Must Sit Below the Photon}

There are two possible directions after the break: begin conformal
transformations, or give the older and easier route to the critical
dimension. We choose the latter, with a warning. The easy argument is not
easy, and it is not the deepest reason. It is nevertheless correct and close
to the historical discovery.

Return to the open bosonic string. Acting once with a transverse creation
operator produces two polarization states. Because there is no longitudinal
third state, the excitation must be a massless vector. In units where one
mode-\(1\) quantum raises \(m^2\) by one,
\begin{equation}
  m_{\mathrm{ground}}^2=-1,
  \qquad
  m_{\mathrm{first}}^2=-1+1=0.
  \label{eq:l05-open-intercept}
\end{equation}
The tachyonic ground state remains a problem of the bosonic theory; it is
acknowledged and set aside. The immediate question is where the negative
unit in \eqref{eq:l05-open-intercept} comes from.

Every string mode is a harmonic oscillator. With \(\hbar=1\), its vacuum
energy is
\begin{equation}
  E_{0,n}=\frac{\omega_n}{2}.
\end{equation}
For the string, \(\omega_n=n\), so one transverse coordinate contributes
formally
\begin{equation}
  E_0^{(1)}
  =\frac12\sum_{n=1}^{\infty}n.
  \label{eq:l05-formal-zero-point}
\end{equation}

\begin{figure}[H]
  \centering
  \includegraphics[width=0.70\textwidth]{lecture_05_figure_11.png}
  \caption{The zero-point sum \(\sum n/2\) written after the mode frequencies
  are recalled.}
  \lectureframe{01:06:18}
  \label{fig:l05-zero-point-board}
\end{figure}

Now the difficulty is sharp. Every term in
\eqref{eq:l05-formal-zero-point} is positive, and the sum diverges. How can
it supply a finite negative intercept?

\section{Subtract the Motion Before Reading the Spectrum}

Before regulating the sum, recall why internal energy was identified with
mass squared. For a particle carrying a large longitudinal momentum
\(P>0\),
\begin{equation}
  E=\sqrt{P^2+m^2}
   =P+\frac{m^2}{2P}
   +O(P^{-3}).
  \label{eq:l05-large-p-energy}
\end{equation}

\begin{figure}[H]
  \centering
  \includegraphics[width=0.70\textwidth]{lecture_05_figure_12.png}
  \caption{The large-\(P\) expansion separating overall momentum from the
  internal mass term.}
  \lectureframe{01:07:35}
  \label{fig:l05-fast-energy-board}
\end{figure}

The first term is the uninteresting overall motion down the \(z\)-axis. It
is conserved and common to all internal states with the same \(P\).
Subtracting it and rescaling removes the kinematic time dilation:
\begin{equation}
  \begin{aligned}
    E-P&=\frac{m^2}{2P}+O(P^{-3}),\\
    2P(E-P)&=m^2+O(P^{-2}).
  \end{aligned}
  \label{eq:l05-internal-energy}
\end{equation}

The atom provides the physical picture. Boost a hydrogen atom. Its total
energy grows, but the energy difference between two internal levels, as
measured in the laboratory time, shrinks because the atom's clock is
dilated. The factor \(1/P\) in \eqref{eq:l05-large-p-energy} is that
suppression. Multiplying by \(2P\) restores the invariant mass-squared
difference.

The raw ground-state energy therefore need not be literally \(-1\). It may
contain a large state-independent piece associated with the overall
longitudinal energy, plus a finite internal remainder. The lecture rewrites
the boosted relation to display both pieces before turning back to the
oscillator sum.

\begin{figure}[H]
  \centering
  \includegraphics[width=0.76\textwidth]{lecture_05_figure_13.png}
  \caption{The invariant mass term displayed beside the large
  state-independent longitudinal contribution.}
  \lectureframe{01:11:34}
  \label{fig:l05-infinite-piece-board}
\end{figure}

Infinite vacuum energies are familiar in quantum field theory. One often
separates a divergent state-independent constant from observable energy
differences. The lecture compares this with the zero-point energy of the
electromagnetic field in the room and openly calls the procedure
unsatisfactory but standard. Here we follow the same heuristic: isolate the
divergent part and ask whether a definite finite remainder survives.

\section{Regulating the Sum}

Introduce a positive regulator \(\epsilon\):
\begin{equation}
  S_1(\epsilon)
  =\sum_{n=1}^{\infty}n e^{-n\epsilon}.
  \label{eq:l05-regulated-sum}
\end{equation}
For every \(\epsilon>0\), the exponential eventually defeats the factor
\(n\), so the sum converges. Sending \(\epsilon\) toward zero removes the
cutoff.

\begin{figure}[H]
  \centering
  \includegraphics[width=0.78\textwidth]{lecture_05_figure_14.png}
  \caption{The exponential regulator applied term by term to
  \(1+2+3+\cdots\).}
  \lectureframe{01:16:42}
  \label{fig:l05-regulator-board}
\end{figure}

The factor \(n\) can be generated by differentiation:
\begin{equation}
  S_1(\epsilon)
  =
  -\frac{d}{d\epsilon}
  \sum_{n=1}^{\infty}e^{-n\epsilon}.
  \label{eq:l05-derivative-trick}
\end{equation}

\begin{figure}[H]
  \centering
  \includegraphics[width=0.76\textwidth]{lecture_05_figure_15.png}
  \caption{Differentiation with respect to the cutoff brings down
  \(-n\).}
  \lectureframe{01:18:05}
  \label{fig:l05-derivative-board}
\end{figure}

The remaining series is geometric:
\begin{equation}
  S_0(\epsilon)
  \equiv
  \sum_{n=1}^{\infty}e^{-n\epsilon}
  =
  \frac{e^{-\epsilon}}{1-e^{-\epsilon}}
  =
  \frac{1}{e^\epsilon-1}.
  \label{eq:l05-geometric-sum}
\end{equation}

\begin{figure}[H]
  \centering
  \includegraphics[width=0.68\textwidth]{lecture_05_figure_16.png}
  \caption{The regulated geometric series in closed form.}
  \lectureframe{01:19:12}
  \label{fig:l05-geometric-board}
\end{figure}

Expand at small \(\epsilon\):
\begin{equation}
  S_0(\epsilon)
  =
  \frac{1}{\epsilon}
  -\frac12
  +\frac{\epsilon}{12}
  -\frac{\epsilon^3}{720}
  +O(\epsilon^5).
  \label{eq:l05-bernoulli-expansion}
\end{equation}
The board obtains this by expanding numerator and denominator separately,
factoring one power of \(\epsilon\), and then expanding the reciprocal. That
calculation contains several in-class sign checks; the compact expression
\eqref{eq:l05-bernoulli-expansion} is its cleaned result.

\begin{figure}[H]
  \centering
  \includegraphics[width=0.78\textwidth]{lecture_05_figure_17.png}
  \caption{The numerator and denominator expansions before the reciprocal is
  simplified.}
  \lectureframe{01:23:10}
  \label{fig:l05-expansion-board}
\end{figure}

Applying \(-d/d\epsilon\) gives
\begin{equation}
  S_1(\epsilon)
  =
  \frac{1}{\epsilon^2}
  -\frac{1}{12}
  +\frac{\epsilon^2}{240}
  +O(\epsilon^4).
  \label{eq:l05-regulated-result}
\end{equation}

\begin{figure}[H]
  \centering
  \includegraphics[width=0.78\textwidth]{lecture_05_figure_18.png}
  \caption{The divergent \(1/\epsilon^2\) term and finite
  \(-1/12\) remainder at the end of the regulator calculation.}
  \lectureframe{01:29:48}
  \label{fig:l05-minus-twelve-board}
\end{figure}

It is common to summarize this result jokingly as
\(1+2+3+\cdots=-1/12\). That is not ordinary convergence. Equation
\eqref{eq:l05-regulated-result} says that this particular regulator produces
a divergent term plus a regulator-independent finite part:
\begin{equation}
  \underbrace{\frac{1}{\epsilon^2}}_{\text{divergent}}
  \;+\;
  \underbrace{\left(-\frac{1}{12}\right)}_{\text{finite}}
  \;+\;O(\epsilon^2).
\end{equation}
The lecture absorbs the first term into the already-present
state-independent piece and retains the second. It also preserves the
warning attributed to Pauli: being infinite does not by itself make a
quantity zero. A consistent theory must explain why the subtraction is
allowed.

Restoring the factor of one-half in
\eqref{eq:l05-formal-zero-point}, one transverse coordinate contributes
\begin{equation}
  E_0^{(1)}=-\frac{1}{24}.
  \label{eq:l05-one-transverse-zero-point}
\end{equation}

\begin{figure}[H]
  \centering
  \includegraphics[width=0.68\textwidth]{lecture_05_figure_19.png}
  \caption{The factor of one-half converts the finite
  \(-1/12\) into \(-1/24\) per transverse coordinate.}
  \lectureframe{01:33:42}
  \label{fig:l05-minus-twenty-four-board}
\end{figure}

\section{Counting the Transverse Directions}

Let \(D_\perp\) be the number of transverse string coordinates. Their total
vacuum contribution is
\begin{equation}
  E_0^{\mathrm{open}}
  =-\frac{D_\perp}{24}.
  \label{eq:l05-open-vacuum}
\end{equation}
For the first open-string oscillator to raise the ground state from
\(-1\) to the massless level, we need
\begin{equation}
  -\frac{D_\perp}{24}=-1.
\end{equation}
Therefore
\begin{equation}
  D_\perp=24.
\end{equation}
Those are the directions perpendicular to the large longitudinal momentum.
Adding that longitudinal spatial direction and time gives
\begin{equation}
  D=D_\perp+2=26.
  \label{eq:l05-critical-dimension}
\end{equation}

\begin{classroomqa}
\audiencequestion{When did this argument enter string
theory? \lecturetimestamp{01:35:27}}
\lecturerresponse{It belongs to the original hadronic-string era, around
1969 and the early 1970s. Regge trajectories were already known
experimentally. Once the oscillator model was explored on its own, the
first excited state was seen to have two rather than three vector
polarizations, forcing it to be massless. The required intercept then led
quickly to twenty-four transverse oscillators.}
\end{classroomqa}

At the time this looked almost too strange to believe. Why should a
regularized vacuum sum know the dimension of spacetime? Yet the pieces fit:
the massless vector, the intercept, and the transverse oscillator count all
selected the same number.

\subsection{The closed string checks the count}

A closed string needs one left-moving and one right-moving mode to reach its
first physical massless level. Its ground state must therefore begin two
units below zero.

\begin{classroomqa}
\audiencequestion{If twenty-four transverse directions give only \(-1\),
must a closed string live in forty-eight transverse
directions? \lecturetimestamp{01:37:47}}
\lecturerresponse{No. The same spacetime contains both open and closed
strings. A closed string already has twice the oscillator content: one tower
circulating in each world-sheet direction. Each tower contributes
\(-24/24=-1\), so together they give \(-2\).}
\end{classroomqa}

Thus
\begin{equation}
  E_0^{\mathrm{closed}}
  =
  2\left(-\frac{D_\perp}{24}\right)
  =-2
  \qquad (D_\perp=24).
\end{equation}
Adding one oscillator from each sector raises the state by two units and
lands at zero mass squared, exactly where the graviton multiplet appeared in
the preceding lecture.

This agreement makes the argument smell better, but it is still not the deep
explanation. In modern quantization the normal-ordering constant, closure of
the Lorentz or Virasoro algebra, and cancellation of the conformal anomaly
belong to one consistent structure. That structure fixes both the intercept
and the critical dimension; the finite cutoff remainder is the historical
way into it.
\editorialnote{The regulator calculation is a heuristic view of the
normal-ordering constant. The deeper derivation of \(D=26\) requires
world-sheet conformal invariance, equivalently cancellation of the conformal
anomaly and consistent constraint algebra.}

That deeper route is next. Only a little complex analysis is needed to begin:
complex coordinates, analytic functions, and the Cauchy--Riemann equations.
The promised physical image is Turkish taffy---pulled one way, stretched
another, reshaped without destroying the structure that matters. The
mathematics behind that image is conformal invariance.
